\chapter{Cross-cutting concepts}\label{ch:crosscutting}
\chapterstatus{draft}{2026-11-16}{\src{CLAUDE.md}, \src{src/server}, \src{src/jobs/worker.ts}}

These are the patterns every change has to respect. Each one exists to protect a quality goal
from \cref{ch:context}.

\section{Notifications through an outbox}\label{sec:outbox}
A change and the messages it causes are written in \emph{one} database transaction:
\code{notify(tx, …)} (\src{src/server/services/notify.ts}) inserts the in-app notification, the
\code{EmailOutbox} rows and the \code{PushMessage} rows next to the change itself. Nothing is sent
from a web request. If the change rolls back, so do its messages.

The worker (\src{src/jobs/worker.ts}) delivers them. It wakes as soon as a transaction that queued
something commits, through PostgreSQL \code{LISTEN}/\code{NOTIFY} (\src{src/server/outbox-signal.ts}),
with a slow poll as a backstop. Each row is claimed with a conditional update, so two workers share
the work instead of duplicating it. Delivery is \emph{at least once}: a crash between sending and
recording sends that one message again. A failed email is retried at most five times with
exponential backoff (1, 2, 4, 8 minutes), then marked \code{FAILED} and shown to admins
(\route{/admin/email}). A deployment can't run without SMTP: it refuses to start
(\src{src/server/startup-checks.ts}), so an email is never just printed; in local development
without SMTP, the worker notes only the recipient and subject, and \route{/dev/mail} shows the
rest.

\paragraph{Delayed and superseded messages.} An outbox row can wait: \code{sendAfter} holds a
student's own booking confirmation (10 minutes) and released marks (5 minutes), and the worker's
backstop poll (60\,s) picks them up when due. A \code{groupKey} ties messages about the same
booking or marks together: \code{notify} deletes unsent emails and unread items with that key in
the same transaction before adding the new one, so people get the current state once
(\cref{sec:rules-usage} has the rules; ADR-018). Once an email has gone it is never recalled, and
the next message says what changed.

\begin{decision}{ADR-007}
  No RDS Proxy or transaction-mode pooler in front of PostgreSQL: \code{LISTEN} needs a direct
  session. Behind a pooler the worker would fall back to polling every \envvar{WORKER_INTERVAL_MS}.
\end{decision}

\section{Live updates by polling}\label{sec:live-polling}
Pages stay current without websockets or server-sent events: API Gateway's HTTP API ends a
request after 30 seconds, and at this scale polling costs cents (\cref{ch:cost}).
\begin{itemize}
  \item \code{<LiveBell>} (\src{src/components/live-bell.tsx}) checks
    \route{/api/notifications/unread} every 30 seconds while the tab is visible, straight away
    when the tab regains focus, and when a push arrives (the service worker posts
    \code{slotty:push} to open tabs). Only the badge changes, except on \route{/notifications},
    \route{/dashboard} and \route{/bookings}, which reload their data when the count goes up.
    A 401 stops the polling.
  \item \code{<LiveRefresh>} reloads a whole page's data every 30 seconds on demo day
    (\route{/today}), and \code{<Countdown>} ticks locally and reloads once its target passes.
\end{itemize}
Both use \code{router.refresh()}, which keeps scroll position and form input.

\section{Access control}
Every service operation receives an \code{Actor} and checks it before reading or writing
(\src{src/server/services/access.ts}): \code{assertCourseRole}, \code{assertStaff},
\code{assertAssignmentRole}, \code{assertAdmin}. Mutations pass \code{{ write: true }} so that
archived courses stay read-only. The role model is described in \cref{ch:journeys}; the rules in
\cref{ch:rules}.

\section{Rules and errors}
Domain functions return a \code{RuleResult} (\src{src/domain/result.ts}); services turn a failed
result into a \code{DomainError} with \code{assertRule}. Server Actions call services through
\code{run()} (\src{src/server/action-utils.ts}), which shows a \code{DomainError}'s message to the
user and hides anything else behind a generic error.

\section{Concurrency}
Booking is the contended path. Capacity is protected by locking the slot rows
(\code{SELECT … FOR UPDATE} in \code{bookings.lockSlots}) inside the booking transaction, and
\enquote{one active booking per student per assignment} by a partial unique index
(\code{booking_one_active_per_student}), so the database refuses a double booking even if two
requests pass the checks at the same moment.

\section{Time}
Times are stored in UTC and shown and entered in the course timezone (\src{src/lib/time.ts}:
\code{fmt}, \code{fmtTimeRange}, \code{fromLocalInput}). Each user picks a 12- or 24-hour clock;
data formats (CSV, form values) stay 24-hour.

\section{Audit}
Moderation-relevant changes are recorded with \code{audit(tx, actor, …)} in the same transaction
as the change (\src{src/server/services/audit.ts}). Admins read them at \route{/admin/audit}.

The log is append-only. A database trigger, \code{AuditLog_append_only}, refuses every
\code{UPDATE} and \code{DELETE} on \code{AuditLog}, from the application, a migration or a
console session alike. A user with audit entries can't be deleted either, because the
foreign key's \code{SET NULL} would change the entries: anonymise the account instead (\#42).
The trigger doesn't cover \code{TRUNCATE} (the tests use it to reset their database) or the
table's owner disabling it; a least-privilege application role, separate from the migration
role, will close those in a later release.

\section{Usage analytics}\label{sec:usage-analytics}
First-party only (ADR-015). In the signed-in layout, \code{<Analytics>} configures the tracker
(\src{src/lib/analytics.ts}) and records page views on navigation, Core Web Vitals through
\code{useReportWebVitals}, uncaught errors and rejections, and app installs; the \route{(app)}
error boundary (\src{src/app/(app)/error.tsx}) reports render errors with their digest.
\code{<ActionForm track="booking">} records a flow's \code{submit} and its \code{done} or
\code{error}; \code{<TrackView>} records its \code{view}. Events queue in memory and are sent to
\route{/api/events} with \code{sendBeacon} every 10 seconds, at 20 events, and when the tab is
hidden or closed; a per-tab random id in \code{sessionStorage} (no cookie) lets funnels follow one
visit. The route handler checks the origin and the session, then \code{recordEvents} validates
every event against the catalog and stores it (\cref{sec:rules-usage}). Reports and feedback are
read only by the platform account at \route{/platform} (ADR-016), whose own browsing is never
recorded. Analytics never breaks
the app: the endpoint always answers 204 and logs its own failures. The worker folds events into
\code{AnalyticsDaily} once an hour (the last 7 full days, idempotently) and deletes expired rows.

Surveys ride on Server Actions: a successful booking, reschedule or mark submission calls
\code{claimSurvey}, and the action's result carries \code{survey}, which \code{<ActionForm>} hands
to \code{<SurveyHost>}. A problem in \code{claimSurvey} is logged and never turns the successful
action into an error.

Production sends a \code{Content-Security-Policy-Report-Only} header (\src{next.config.ts}):
browsers report, without blocking, anything outside \code{'self'}. The first finding is an
\code{eval} probe (zod's JIT check); enforcing the policy needs zod's \code{jitless} setting and
script nonces first.

\section{Configuration and start-up}
All configuration is read from the environment at run time (\cref{app:config}), so one image runs
locally and on AWS. Production refuses to start with configuration that would silently
break links (\src{src/server/startup-checks.ts}).

\section{Requests from the edge}
On AWS, requests arrive through CloudFront and API Gateway. \src{src/proxy.ts} (Next.js Proxy)
refuses any request without CloudFront's secret \code{x-origin-verify} header and resets the
forwarded host and protocol to \envvar{APP_URL}, because API Gateway replaces them with its own
hostname and Next's Server Actions compare the forwarded host with the \code{Origin} header
(\cref{sec:aws-request-path}). With \envvar{ORIGIN_VERIFY_SECRET} unset (development) it does
nothing. \route{/health/live} answers without touching the database, for container health
checks; \route{/health} also checks the database and reports the running release
(\envvar{APP_VERSION}), for monitoring and for the deploy's last check (\cref{sec:deploy}).

\begin{todo}
  \begin{itemize}
    \item Authentication and sessions in detail: password hashing (Argon2), session cookies and
      their flags, token lifetimes, sign-out behaviour.
    \item Caching and revalidation in the App Router; what the PWA service worker caches.
    \item Logging conventions and what must never be logged (tokens, passwords, marks).
  \end{itemize}
\end{todo}
